%!TEX root = main.tex
\section{Background and Related Work}
\label{sec:background}

\subsection{OpenEdgeCGRA Architecture}
\label{subsec:cgra-basics}

Coarse Grained Reconfigurable Arrays (CGRAs) are primarily composed of an array of interconnected
PEs (or Reconfigurable Cells --RCs--) that can be configured dynamically to implement a variety of computational tasks. The PEs are connected through a fixed interconnect, executing in lockstep under a configuration bitstream loaded before kernel execution~\cite{cgra-survey}. Each PE contains a word-width ALU, a small register file, and neighbor connections. Unlike FPGAs (bit-level reconfiguration, high overhead), CGRAs reconfigure at the instruction level with microsecond latency.

OpenEdgeCGRA~\cite{openedge-cgra} provides: \textit{(i)} a $4\times4$ toroidal mesh of PEs (i.e 16 PEs),  where column~3 wraps to column~0 and row~3 to row~0, giving maximum distance of 4 hops between any two PEs; \textit{(ii)} per-PE resources including a 32-bit ALU (1~clock cycle), $32\!\times\!32$ multiplier (3~clock cycles), four registers (R0--R3), output register ROUT visible to neighbors, and branchless select instructions (BSFA/BZFA); \textit{(iii)} shared scratchpad SRAM with 2-cycle latency (one access/column/cycle) organized as 4 interleaved banks; and \textit{(iv)} tight coupling with a RISC-V host core (X-HEEP~\cite{xheep}). At SoC level, we consider two host-to-memory organizations. The simpler one-to-many path (\textit{OneToM}) forces host and accelerator traffic to share a common route to memory, whereas the richer \textit{N:M} crossbar sustains more parallel access to the interleaved SRAM banks. As a result, \textit{OneToM} incurs more serialization and contention, while \textit{N:M} reduces memory-induced stalls.

\textbf{Configuration and execution model.} A kernel is encoded as a sequence of 32-bit \emph{configuration words}---one per PE per cycle---stored in a Configuration Register File (CRF) with $D$ entries per column---32 entries in our baseline CGRA. Before dispatch, the host CPU writes configuration words into an SRAM backing store called the \emph{Context Memory} (CMEM) via memory-mapped I/O (MMIO). During execution, each column's program counter advances through CRF entries in lockstep, dispatching one instruction per PE. When a kernel exceeds the CRF capacity, the host must reload the CMEM with the next batch of instructions before the CGRA can resume---a transparent but costly operation whose implications are analyzed in Section~\ref{subsec:creg-analysis}.

\subsection{CGRA Programming: Limitations for ZKP Workloads}
\label{subsec:cgra-programming}

\textbf{Manual scheduling} requires placing instructions among $NB$ bundles of PEs slots,
satisfying data dependencies, multi-cycle latencies, memory-port arbitration, and register pressure---producing CSV files of $\sim$4{,}750 lines for our main kernel, with errors detectable only through simulation.

\noindent \textbf{Automatic compilers} (Compigra~\cite{compigra}, Wang et al.~\cite{wang2025mlir}) accept C code and use heuristic placement, working well for regular DSP kernels but encountering structural limitations for cryptographic modular arithmetic (our use case): (1)~multi-precision limb decomposition is invisible at the C level; (2)~toroidal routing patterns have no C-level representation; (3)~architectural primitives like BSFA (\textit{Branch/Select on Flag A}, a branchless correction needed in our main kernel) are inaccessible; and (4)~heuristic scheduling forfeits cycle-count predictability needed for real-time ZKP functions.

\subsection{Poseidon2 S-box Exponentiation}
\label{subsec:poseidon2}

Poseidon2~\cite{poseidon2} is a ZKP-friendly hash function engineered to minimize finite-field multiplications per permutation, since each multiplication translates directly to proof-generation cost. Fig.~\ref{fig:poseidon2-structure} shows its round structure: it operates on a state $s = (s_0, s_1, s_2, \ldots, s_{t-1})$ of $t$ field elements through $R_F/2$ full rounds, then a $R_P$ partial round for cell $s_0$, and next $R_F/2$ full rounds again on $s$. Each round applies a S-box layer ($x \mapsto x^d \Mod p$), a linear {\textit Maximum Distance Separable} layer (MDS, a $t \times t$ matrix used to provide maximum diffusion), and a round-constant addition ($+c_i$, a standard feature to ensure each round behaves differently).
For a typical parameterization ($t\!=\!16$, $R_F\!=\!8$, $R_P\!=\!13$), Poseidon2 invokes $8 \times 16 + 13 = 141$ S-box evaluations, making the S-box layer over 90\% of total computation on a 32-bit CGRA.

\begin{figure}[t]
\centering
\resizebox{\columnwidth}{!}{%
\begin{tikzpicture}[
    % 1. Configuración de la cadena para el posicionamiento automático
    start chain=going right,
    node distance=0.12cm, 
    % Estilos
    box/.style={draw, thick, inner sep=3pt, minimum height=0.55cm, font=\footnotesize\sffamily, rounded corners=2pt, on chain},
    sbox/.style={box, fill=cgrared!15},
    mds/.style={box, fill=cgrablue!15},
    rc/.style={box, fill=cgragreen!15},
    arr/.style={-{Stealth[length=4pt]}, thick},
    lbl/.style={font=\scriptsize\sffamily, text=cgragray},
    group_label/.style={font=\footnotesize\sffamily\bfseries}
]

% --- DEFINICIÓN ESPACIAL (Nodos) ---
\node[on chain, font=\footnotesize\sffamily, inner sep=2pt] (inp) {state};

\node[sbox] (sb1) {S-box};
\node[mds]  (m1)  {MDS};
\node[rc]   (r1)  {$+c_i$};

\node[sbox] (sb2) {S-box};
\node[mds]  (m2)  {MDS};
\node[rc]   (r2)  {$+c_i$};

\node[sbox] (sb3) {S-box};
\node[mds]  (m3)  {MDS};

\node[on chain, font=\footnotesize\sffamily, inner sep=2pt] (out) {out};


% --- ENRUTADO (Conexiones) ---
\draw[arr] (inp) -- (sb1);
\draw[arr] (sb1) -- (m1);
\draw[arr] (m1)  -- (r1);

\draw[arr] (r1)  -- node[above, lbl] {$\cdots$} (sb2);
\draw[arr] (sb2) -- (m2);
\draw[arr] (m2)  -- (r2);

\draw[arr] (r2)  -- node[above, lbl] {$\cdots$} (sb3);
\draw[arr] (sb3) -- (m3);
\draw[arr] (m3)  -- (out);


% --- ETIQUETADO ---
\node[lbl, above=0.05cm of sb1] {$\forall s_i$};
\node[group_label, above=0.35cm of m1] {Full ($\times R_F/2$)};

\node[lbl, above=0.05cm of sb2] {$s_0$};
\node[group_label, above=0.35cm of m2] {Partial ($\times R_P$)};

\node[lbl, above=0.05cm of sb3] {$\forall s_i$};
\node[group_label, above=0.35cm of m3] {Full ($\times R_F/2$)};

\end{tikzpicture}%
}
\caption{Poseidon2 permutation structure. The S-box ($x^d \Mod p$) dominates cost.}
\label{fig:poseidon2-structure}
\end{figure}

\subsection{Related Work}
\label{subsec:related-hw}

Table~\ref{tab:related-work} positions our contribution relative to existing ZKP hardware accelerators.
PipeZK~\cite{pipezk} proposes an ASIC pipeline for MSM (\textit{Multi-Scalar Multiplication}) and NTT (\textit{Number Theoretic Transform}), which
are the two most time-consuming computational primitives for generating Zero-Knowledge Proofs (zk-SNARKs), but it does not address ZK-friendly hashing. AMAZE~\cite{amaze} implements MiMC (a ZK-friendly hash function) on FPGA with fixed RTL, achieving $13\times$ over CPU. HashEmAll~\cite{hashemall} provides FPGA-optimized cores for Rescue-Prime, Griffin, and Reinforced Concrete (up to $23\times$ over CPU on BN254), yet it targets different hash families than Poseidon2. TRIDENT~\cite{trident} accelerates the original Poseidon~\cite{poseidon} on FPGA for BLS12-381 (Filecoin).

Although FPGAs are inherently reprogrammable platforms, the accelerators listed above implement their arithmetic pipelines as fixed RTL designs optimized for a specific prime field and hash function. Adapting them to a different field (e.g., from BLS12-381 to BabyBear) or a different hash (e.g., from MiMC to Poseidon2) would require modifying the HDL source, re-running synthesis and place-and-route, and generating a new bitstream---a process that takes minutes to hours and demands hardware design expertise. A CGRA operates at a different abstraction level: the hardware (ALUs, registers, interconnect) remains fixed, and switching to a new kernel amounts to loading a new instruction configuration via DMA, completing in microseconds with no hardware re-synthesis. This distinction is particularly relevant for wearable ZKP devices, where the proof system may evolve (new field, new exponent, new hash) and the device must adapt through a software update rather than FPGA re-flashing.

Separately, on the CGRA programming side, efficient CGRA compilers must support both spatially configured and time-multiplexed tasks. Compigra~\cite{compigra} and Wang et al.~\cite{wang2025mlir} provide automatic C-to-CGRA compilation tools that represent the state-of-the-art and consider some of the mentioned spatial-temporal challenges. However, they target general-purpose kernels and lack mechanisms for multi-precision spatial arithmetic, toroidal routing, or branchless correction---features essential for ZKP workloads.

To our knowledge, no prior work has mapped Poseidon2 hash function onto a CGRA, designed an API for spatial-temporal cryptographic programming, or analyzed Configuration Register File capacity pressure for ZKP-class workloads.

\begin{table}[t]
\caption{Comparison with related ZKP hardware accelerators.
The ``Reprog.'' column reflects the cost of switching to a different hash or field.}
\label{tab:related-work}
\centering
\scriptsize
\setlength{\tabcolsep}{2pt}
\begin{tabular}{@{}lllll@{}}
\toprule
\textbf{Work} & \textbf{Platf.} & \textbf{Primitive} & \textbf{Field} & \textbf{Reprog.} \\
\midrule
PipeZK~\cite{pipezk} & ASIC & MSM+NTT & BN254 & Fabr. \\
AMAZE~\cite{amaze} & FPGA & MiMC & BN254 & RTL \\
HashEmAll~\cite{hashemall} & FPGA & Resc./Griff. & BN254 & RTL \\
TRIDENT~\cite{trident} & FPGA & Poseidon & BLS12 & RTL \\
\textbf{This work} & \textbf{CGRA} & \textbf{Poseidon2} & \textbf{BabyB.} & \textbf{DMA} \\
\bottomrule
\end{tabular}
\end{table}
